\ExampleChapter{示例论文使用说明}{示例论文使用说明}
本工程是一份可编译、可修改、可追溯的 GZUThesis 排版示例。
正文采用无人机飞行数据异常检测这一主题组织，但不构成可以直接提交的学位论文。
其中，文献条目对应真实出版物；数据、飞行编号和故障过程由本工程脚本合成。
所有实验表格均由该脚本计算，未填入其他论文的测试成绩，也未把演示数值标记为真实飞行结论。

示例覆盖封面、中英文摘要、目录、插图清单、附表清单、符号说明、六章正文、
致谢、参考文献和三个附录。正文同时展示分层标题、双语题注、并排子图、
三线表、分组表头、跨页长表、公式、矩阵、定理、证明、算法和代码清单。
这些内容的主要用途是检验模板接口和观察版式，而不是扩充论文字数。

默认入口为 \nolinkurl{main.tex}，使用 BibLaTeX 与 Biber。
工程还提供 BibTeX 入口和两种著录格式的切换入口。
各入口共用同一份正文和 \nolinkurl{ref/references.bib}，因此切换文献系统时不必复制章节。
编译命令与依赖说明见工程根目录的 \nolinkurl{README.md}。

当前模板尚未提供正式的摘要环境和关键词配置，本示例在
\nolinkurl{config/commands.tex} 中使用带 \texttt{Example} 前缀的辅助命令生成这些前置页面。
这不等于模板已完成对应功能。答辩委员会名单、原创性声明和授权页没有伪造，
正式使用时应根据学校提供的原始材料另行制作。

\section*{阅读与修改建议}
首先检查封面信息和版面；然后依次查看图表、数学和参考文献测试；
最后查看附录中的配置字典及可复现清单。
将本工程复制到独立目录后再修改，避免覆盖正在维护的模板源码。
\nolinkurl{gzuthesis.dtx} 是随包固定的类文件源代码快照，
不是要求用户以后在每篇论文中重新编辑的排版参数文件。

\section*{数值与统计边界}
本例包含 12 个合成架次，每架次 1000 个采样点，采样率为
\qty{20}{\hertz}。训练、校准和测试按架次划分，分别使用 6、3、3 个架次。
测试异常仅用于计算评价指标，不参与模型拟合和阈值选择。
本例不使用点调整，也不宣称任何置信区间或统计显著性。
